\subsection{Backpropagation Through Time (BPTT)}

\definition{Backpropagation Through Time (BPTT)} \textbf{is the algorithm used to train Recurrent Neural Networks (RNNs)}. It is simply \textbf{backpropagation} (\autopageref{sec:backpropagation-conceptual-introduction}) \textbf{applied to the unrolled RNN}. However, an RNN is \textbf{not a standard feed-forward network} because it has \textbf{cycles} due to its recurrent connections:
\begin{equation*}
    h_{t-1} \rightarrow \text{[RNN cell]} \rightarrow h_{t} \rightarrow \text{[RNN cell]} \rightarrow h_{t+1} \rightarrow \ldots
\end{equation*}
These cycles create \textbf{dependencies over time}, making the RNN \textbf{deep in time} rather than deep in space. So, the RNN cannot run backpropagation directly on its cyclic graph. The solution is to \textbf{unroll the RNN over time}, converting it into a \textbf{Directed Acyclic Graph (DAG)}. Once unrolled, we can perform backpropagation as usual, computing gradients for each time step and accumulating them to update the shared weights.

\highspace
Once unrolled, we simply:
\begin{enumerate}
    \item Run forward pass through all timesteps:
    \begin{equation*}
        \text{compute } h_1, \, h_2, \, \ldots, \, h_T
    \end{equation*}
    \item Compute the loss at each timestep and sum them:
    \begin{equation*}
        E = \sum_{t=1}^{T} E_t
    \end{equation*}
    \item Run backpropagation through the unrolled network, computing gradients at each timestep and accumulating them:
    \begin{equation*}
        \frac{\partial E}{\partial W_x} = \sum_{t=1}^{T} \frac{\partial E_t}{\partial W_x}, \quad
        \frac{\partial E}{\partial W_h} = \sum_{t=1}^{T} \frac{\partial E_t}{\partial W_h}, \quad
        \frac{\partial E}{\partial b} = \sum_{t=1}^{T} \frac{\partial E_t}{\partial b}
    \end{equation*}
    \item Update the shared weights using the accumulated gradients.
\end{enumerate}
That is Backpropagation Through Time in a nutshell. It is just plain backpropagation but applied \textbf{through time}, not space. But, \textbf{how do we unroll an RNN?} Let's explore that next.

\newpage

\begin{deepeningbox}[: Why backpropagation cannot be applied directly to RNNs?]
    In a Recurrent Neural Network, the main issue is \textbf{not} that a parameter is used multiple times, but that the computational graph contains \textbf{cycles}.

    \highspace
    In a standard feed-forward neural network, the computation flows in one direction:
    \begin{equation*}
        \text{input } x \rightarrow \text{layer } W \rightarrow \text{hidden } h \rightarrow \text{loss } E
    \end{equation*}
    This defines a \textbf{Directed Acyclic Graph (DAG)}, which is required for backpropagation to work.

    \highspace
    In contrast, an RNN introduces a \textbf{recurrent connection}:
    \begin{equation*}
        h_t = f(W_x x_t + W_h h_{t-1} + b)
    \end{equation*}
    This creates a \textbf{cycle}:
    \begin{equation*}
        h_{t-1} \rightarrow h_t \rightarrow h_{t+1} \rightarrow \ldots
    \end{equation*}

    \textcolor{Red2}{\faIcon{exclamation-triangle} \textbf{Why is this a problem?}}
    Because backpropagation requires a DAG, but the RNN defines a \textbf{cyclic (recursive) computation}. Therefore, we cannot directly apply backpropagation.

    \highspace
    \textcolor{Red2}{\faIcon{question-circle} \textbf{Let's see from a mathematical perspective.}} In an RNN:
    \begin{equation*}
        h_t = f(W_x x_t + W_h h_{t-1} + b)
    \end{equation*}
    The term $h_{t-1}$ is a function of $h_{t-2}$, which is a function of $h_{t-3}$, and so on. This is a \textbf{nested function composition}:
    \begin{equation*}
        h_t = f\left(W_x x_t + W_h f\left(W_x x_{t-1} + W_h f\left(W_x x_{t-2} + \ldots\right)\right)\right)
    \end{equation*}
    Now, if we try to compute the gradient $\frac{\partial E}{\partial W_h}$, because we want to update $W_h$, we would have to apply the chain rule through this nested composition. This means we would have to compute:
    \begin{equation*}
        \frac{\partial E}{\partial W_h} = \frac{\partial E}{\partial h_t} \cdot \frac{\partial h_t}{\partial W_h}
    \end{equation*}
    But $\frac{\partial h_t}{\partial W_h}$ depends on:
    \begin{equation*}
        \frac{\partial h_t}{\partial W_h} = \frac{\partial f}{\partial W_h} + \frac{\partial f}{\partial h_{t-1}} \cdot \frac{\partial h_{t-1}}{\partial W_h}
    \end{equation*}
    And $\frac{\partial h_{t-1}}{\partial W_h}$ depends on:
    \begin{equation*}
        \frac{\partial h_{t-1}}{\partial W_h} = \frac{\partial f}{\partial W_h} + \frac{\partial f}{\partial h_{t-2}} \cdot \frac{\partial h_{t-2}}{\partial W_h}
    \end{equation*}
    And so on, leading to a \hl{recursive definition of the gradient, where each term depends on the previous one.} To compute it explicitly, we \hl{must expand this recursion over a finite number of time steps.}

    \highspace
    \textcolor{Green3}{\faIcon{check-circle} \textbf{Solution: BPTT (unrolling).}}
    We convert the cyclic graph into a \textbf{finite, acyclic graph} by unrolling the RNN over $T$ time steps:
    \begin{equation*}
        x_1 \rightarrow h_1 \rightarrow h_2 \rightarrow \ldots \rightarrow h_T
    \end{equation*}
    Now the model becomes equivalent to a deep feed-forward network, and backpropagation can be applied.

    \highspace
    \textcolor{DarkOrange3}{\faIcon{question-circle} \textbf{What changes after unrolling?}}
    Although unrolling enables backpropagation, it reveals that the same weights are used at every time step. Therefore, each parameter influences the loss through \textbf{many paths across time}:
    \begin{equation*}
        \frac{\partial E}{\partial W_h} = \sum_{t=1}^{T} \frac{\partial E}{\partial h_t} \cdot \frac{\partial h_t}{\partial W_h}
    \end{equation*}
    This leads to \textbf{long chains of dependencies}, which in turn cause the well-known \textbf{vanishing and exploding gradient problems}. But we will discuss that separately.

    \highspace
    \textcolor{Green3}{\faIcon{question-circle} \textbf{But why not use all time steps?}} 
    In principle, the correct gradient requires considering \textbf{all previous time steps}. However, in practice this is often inefficient:
    \begin{itemize}
        \item The computational cost grows linearly with the sequence length $T$, requiring both \textbf{more memory} (to store all intermediate states) and \textbf{more time} for backpropagation.
        \item More importantly, contributions from distant time steps tend to become \textbf{negligible} due to the repeated multiplication of Jacobians, leading to the \textbf{vanishing/exploding gradient problem}.
    \end{itemize}
    We will discuss these issues in more detail in the next section, but the key point is that while BPTT is the correct algorithm for training RNNs, it is often impractical to apply it over long sequences due to computational and optimization challenges.

    \highspace
    For these reasons, a common practical solution is \definition{Truncated Backpropagation Through Time (TBPTT)}, where gradients are propagated only for a limited number of steps (i.e., a BPTT but truncated to a fixed window).

    \highspace
    In this course, however, we focus on the \textbf{full BPTT formulation}, as it clearly exposes the limitations of RNNs and motivates more advanced architectures such as \textbf{LSTMs and GRUs}.
\end{deepeningbox}